\section{Audit and completeness}
\begin{table}[t]
\centering
\caption{Result completeness. Smoke records are unplanned diagnostics and are excluded.}
\label{tab:completeness}
\small
\begin{tabular}{lrr}
\toprule
Phase & Completed & Expected \\
\midrule
Part 1 tuning & 270 & 270 \\
Part 1 final & 837 & 837 \\
Bonus tuning & 36 & 36 \\
Bonus final & 108 & 108 \\
\midrule
Total & 1,251 & 1,251 \\
\bottomrule
\end{tabular}
\end{table}


The audit outcome is \texttt{complete\_with\_declared\_exceptions}. All planned runs are complete. The declared exception concerns only the original derived bonus TF-IDF reference; the reporting rebuild fixes the selection rule from intact Part~1 atomic records.

\section{Additional result checks}
\begin{table*}[t]
\centering
\caption{Bonus target-domain macro-F1 (percentage points), mean $\pm$ sample SD ($n=3$). The TF-IDF and target-only Transformer rows reuse Part 1 runs.}
\label{tab:bonus-performance}
\scriptsize
\setlength{\tabcolsep}{3.5pt}
\begin{tabular}{llrrrr}
\toprule
Target & Method & 200 & 500 & 1,000 & 2,000 \\
\midrule
NSMC & Best TF-IDF & 60.62 $\pm$ 0.96 (3) & 63.99 $\pm$ 0.58 (3) & 67.02 $\pm$ 0.28 (3) & 69.65 $\pm$ 0.19 (3) \\
 & Target-only Transformer & 50.77 $\pm$ 3.52 (3) & 51.26 $\pm$ 6.13 (3) & 54.72 $\pm$ 5.97 (3) & 57.90 $\pm$ 1.17 (3) \\
 & Pooled CE & 62.82 $\pm$ 0.83 (3) & 64.73 $\pm$ 1.33 (3) & 66.55 $\pm$ 1.07 (3) & 68.71 $\pm$ 0.38 (3) \\
 & SupCon & 63.08 $\pm$ 0.79 (3) & 64.83 $\pm$ 1.32 (3) & 66.02 $\pm$ 1.80 (3) & 68.52 $\pm$ 0.68 (3) \\
 & DB-CDSupCon & 63.59 $\pm$ 0.75 (3) & 65.36 $\pm$ 0.34 (3) & 66.53 $\pm$ 0.74 (3) & 68.78 $\pm$ 0.62 (3) \\
\addlinespace
Shopping & Best TF-IDF & 71.93 $\pm$ 0.96 (3) & 76.73 $\pm$ 1.52 (3) & 79.81 $\pm$ 0.44 (3) & 81.86 $\pm$ 0.31 (3) \\
 & Target-only Transformer & 63.92 $\pm$ 2.74 (3) & 64.70 $\pm$ 5.53 (3) & 72.51 $\pm$ 2.31 (3) & 71.75 $\pm$ 7.66 (3) \\
 & Pooled CE & 69.11 $\pm$ 3.29 (3) & 74.02 $\pm$ 0.67 (3) & 75.86 $\pm$ 1.58 (3) & 75.45 $\pm$ 2.19 (3) \\
 & SupCon & 69.48 $\pm$ 3.63 (3) & 74.22 $\pm$ 0.85 (3) & 75.52 $\pm$ 1.70 (3) & 78.04 $\pm$ 0.70 (3) \\
 & DB-CDSupCon & 69.73 $\pm$ 3.77 (3) & 74.30 $\pm$ 0.46 (3) & 75.35 $\pm$ 1.70 (3) & 77.88 $\pm$ 0.69 (3) \\
\addlinespace
Steam & Best TF-IDF & 60.50 $\pm$ 0.60 (3) & 62.42 $\pm$ 0.58 (3) & 64.49 $\pm$ 0.46 (3) & 67.32 $\pm$ 0.44 (3) \\
 & Target-only Transformer & 49.18 $\pm$ 2.60 (3) & 54.20 $\pm$ 1.14 (3) & 58.02 $\pm$ 1.87 (3) & 57.33 $\pm$ 5.48 (3) \\
 & Pooled CE & 62.71 $\pm$ 1.04 (3) & 62.83 $\pm$ 1.33 (3) & 65.18 $\pm$ 0.35 (3) & 66.29 $\pm$ 0.41 (3) \\
 & SupCon & 61.70 $\pm$ 0.20 (3) & 63.93 $\pm$ 0.14 (3) & 64.67 $\pm$ 0.57 (3) & 66.51 $\pm$ 0.97 (3) \\
 & DB-CDSupCon & 62.78 $\pm$ 1.50 (3) & 62.83 $\pm$ 0.92 (3) & 64.97 $\pm$ 0.84 (3) & 66.14 $\pm$ 0.66 (3) \\
\addlinespace
\bottomrule
\end{tabular}
\end{table*}

\begin{table}[t]
\centering
\caption{Bonus label-efficiency result. Crossovers compare each method with the fixed best Part 1 TF-IDF reference at the same target size. Savings are reported only if both the method and target-only Transformer crossovers are observed.}
\label{tab:bonus-labels}
\scriptsize
\resizebox{\columnwidth}{!}{%
\begin{tabular}{@{}llrr@{}}
\toprule
Target & Method & Crossover & Labels saved \\
\midrule
NSMC & Target-only Transformer & Not observed & Not estimable \\
 & Pooled CE & 200 & Not estimable \\
 & SupCon & 200 & Not estimable \\
 & DB-CDSupCon & 200 & Not estimable \\
\addlinespace
Shopping & Target-only Transformer & Not observed & Not estimable \\
 & Pooled CE & Not observed & Not estimable \\
 & SupCon & Not observed & Not estimable \\
 & DB-CDSupCon & Not observed & Not estimable \\
\addlinespace
Steam & Target-only Transformer & Not observed & Not estimable \\
 & Pooled CE & 200 & Not estimable \\
 & SupCon & 200 & Not estimable \\
 & DB-CDSupCon & 200 & Not estimable \\
\addlinespace
\bottomrule
\end{tabular}
}
\end{table}

\begin{table}[t]
\centering
\caption{Bonus source retention and representation separation, averaged over the four target-label sizes. Retention is the source macro-F1 difference from pooled CE.}
\label{tab:bonus-diagnostics}
\scriptsize
\begin{tabular}{llrr}
\toprule
Target & Method & Source $\Delta$ & Separation gap \\
\midrule
NSMC & Pooled CE & +0.000 & 0.095 \\
 & SupCon & -0.000 & 0.094 \\
 & DB-CDSupCon & +0.005 & 0.088 \\
\addlinespace
Shopping & Pooled CE & +0.000 & 0.078 \\
 & SupCon & +0.015 & 0.091 \\
 & DB-CDSupCon & +0.015 & 0.086 \\
\addlinespace
Steam & Pooled CE & +0.000 & 0.154 \\
 & SupCon & -0.001 & 0.146 \\
 & DB-CDSupCon & -0.001 & 0.147 \\
\addlinespace
\bottomrule
\end{tabular}
\end{table}


\begin{figure}[ht]
  \centering
  \includegraphics[width=0.62\textwidth]{confusion_comparison.pdf}
  \caption{NSMC full-data confusion matrices summed over three seeds and normalized within the actual class. Parentheses give summed counts.}
  \label{fig:confusion}
\end{figure}

\section{Frozen configurations}
\begin{longtable}{lll >{\raggedright\arraybackslash}p{7.2cm}}
\caption{Frozen Part 1 configurations selected at 8,000 labels with tuning seed 777.}\label{tab:selected-configs}\\
\toprule
Dataset & Model & ID & Configuration \\
\midrule
\endfirsthead
\toprule
Dataset & Model & ID & Configuration \\
\midrule
\endhead
NSMC & TF-IDF + LR & lr06 & \texttt{\{"C": 2.0, "id": "lr06", "max\_iter": 3000, "solver": "liblinear"\}} \\
NSMC & TF-IDF + SVM & svm05 & \texttt{\{"C": 0.25, "id": "svm05"\}} \\
NSMC & TF-IDF + RF & rf09 & \texttt{\{"id": "rf09", "max\_depth": null, "max\_features": "sqrt", "min\_samples\_leaf": 2, "n\_estimators": 700\}} \\
NSMC & TF-IDF + XGBoost & xgb03 & \texttt{\{"colsample\_bytree": 0.8, "id": "xgb03", "learning\_rate": 0.05, "max\_depth": 4, "n\_estimators": 400, "reg\_lambda": 1.0, "subsample": 0.8\}} \\
NSMC & Avg. embeddings + MLP & mlp08 & \texttt{\{"dropout": 0.3, "embedding\_dim": 128, "hidden\_dim": 288, "id": "mlp08", "learning\_rate": 0.002, "num\_layers": 2\}} \\
NSMC & LSTM & lstm08 & \texttt{\{"dropout": 0.3, "embedding\_dim": 128, "hidden\_dim": 112, "id": "lstm08", "learning\_rate": 0.002, "num\_layers": 2\}} \\
NSMC & GRU & gru08 & \texttt{\{"dropout": 0.3, "embedding\_dim": 128, "hidden\_dim": 128, "id": "gru08", "learning\_rate": 0.002, "num\_layers": 2\}} \\
NSMC & Transformer & tr07 & \texttt{\{"dropout": 0.4, "embedding\_dim": 128, "ff\_multiplier": 2, "hidden\_dim": 128, "id": "tr07", "learning\_rate": 0.001, "num\_heads": 4, "num\_layers": 2\}} \\
NSMC & Mamba & mb08 & \texttt{\{"d\_conv": 4, "d\_state": 16, "dropout": 0.3, "embedding\_dim": 128, "expand": 2, "hidden\_dim": 128, "id": "mb08", "learning\_rate": 0.002, "num\_layers": 2\}} \\
Shopping & TF-IDF + LR & lr06 & \texttt{\{"C": 2.0, "id": "lr06", "max\_iter": 3000, "solver": "liblinear"\}} \\
Shopping & TF-IDF + SVM & svm05 & \texttt{\{"C": 0.25, "id": "svm05"\}} \\
Shopping & TF-IDF + RF & rf09 & \texttt{\{"id": "rf09", "max\_depth": null, "max\_features": "sqrt", "min\_samples\_leaf": 2, "n\_estimators": 700\}} \\
Shopping & TF-IDF + XGBoost & xgb03 & \texttt{\{"colsample\_bytree": 0.8, "id": "xgb03", "learning\_rate": 0.05, "max\_depth": 4, "n\_estimators": 400, "reg\_lambda": 1.0, "subsample": 0.8\}} \\
Shopping & Avg. embeddings + MLP & mlp08 & \texttt{\{"dropout": 0.3, "embedding\_dim": 128, "hidden\_dim": 288, "id": "mlp08", "learning\_rate": 0.002, "num\_layers": 2\}} \\
Shopping & LSTM & lstm04 & \texttt{\{"dropout": 0.2, "embedding\_dim": 96, "hidden\_dim": 96, "id": "lstm04", "learning\_rate": 0.002, "num\_layers": 1\}} \\
Shopping & GRU & gru08 & \texttt{\{"dropout": 0.3, "embedding\_dim": 128, "hidden\_dim": 128, "id": "gru08", "learning\_rate": 0.002, "num\_layers": 2\}} \\
Shopping & Transformer & tr07 & \texttt{\{"dropout": 0.4, "embedding\_dim": 128, "ff\_multiplier": 2, "hidden\_dim": 128, "id": "tr07", "learning\_rate": 0.001, "num\_heads": 4, "num\_layers": 2\}} \\
Shopping & Mamba & mb08 & \texttt{\{"d\_conv": 4, "d\_state": 16, "dropout": 0.3, "embedding\_dim": 128, "expand": 2, "hidden\_dim": 128, "id": "mb08", "learning\_rate": 0.002, "num\_layers": 2\}} \\
Steam & TF-IDF + LR & lr06 & \texttt{\{"C": 2.0, "id": "lr06", "max\_iter": 3000, "solver": "liblinear"\}} \\
Steam & TF-IDF + SVM & svm05 & \texttt{\{"C": 0.25, "id": "svm05"\}} \\
Steam & TF-IDF + RF & rf09 & \texttt{\{"id": "rf09", "max\_depth": null, "max\_features": "sqrt", "min\_samples\_leaf": 2, "n\_estimators": 700\}} \\
Steam & TF-IDF + XGBoost & xgb03 & \texttt{\{"colsample\_bytree": 0.8, "id": "xgb03", "learning\_rate": 0.05, "max\_depth": 4, "n\_estimators": 400, "reg\_lambda": 1.0, "subsample": 0.8\}} \\
Steam & Avg. embeddings + MLP & mlp08 & \texttt{\{"dropout": 0.3, "embedding\_dim": 128, "hidden\_dim": 288, "id": "mlp08", "learning\_rate": 0.002, "num\_layers": 2\}} \\
Steam & LSTM & lstm08 & \texttt{\{"dropout": 0.3, "embedding\_dim": 128, "hidden\_dim": 112, "id": "lstm08", "learning\_rate": 0.002, "num\_layers": 2\}} \\
Steam & GRU & gru08 & \texttt{\{"dropout": 0.3, "embedding\_dim": 128, "hidden\_dim": 128, "id": "gru08", "learning\_rate": 0.002, "num\_layers": 2\}} \\
Steam & Transformer & tr08 & \texttt{\{"dropout": 0.3, "embedding\_dim": 128, "ff\_multiplier": 2, "hidden\_dim": 128, "id": "tr08", "learning\_rate": 0.002, "num\_heads": 8, "num\_layers": 2\}} \\
Steam & Mamba & mb07 & \texttt{\{"d\_conv": 3, "d\_state": 16, "dropout": 0.4, "embedding\_dim": 128, "expand": 2, "hidden\_dim": 128, "id": "mb07", "learning\_rate": 0.001, "num\_layers": 2\}} \\
\bottomrule
\end{longtable}


\section{Additional interpretation details}

The machine-readable appendix is more exhaustive than a printed table. \texttt{part1\_metrics\_by\_seed.csv} contains all 837 Part~1 final rows. \texttt{part1\_summary.csv} contains every learning-curve cell, gains from the previous and approximate half-size budgets, and normalized log-size AULC. \texttt{bonus\_metrics\_by\_seed\_corrected.csv} contains all five bonus comparison methods with corrected TF-IDF reuse. \texttt{bonus\_paired\_differences.csv} records every ordered paired seed difference. Cost, source-retention, representation, and label-efficiency files sit beside them.

\paragraph{Loss and deployment.} The SupCon projection is a linear layer, ReLU, second linear layer, and L2 normalization. It adds \Result{bonus_projection_parameters} trainable parameters during SupCon training. The deployed classifier bypasses this head, leaving \Result{bonus_deployed_parameters} parameters for pooled CE, SupCon, and DB-CDSupCon. Thus contrastive training changes fit time and peak training memory but not the deployed computation graph.
